\SourceLeaf{229}%
\bigskip\Needspace{9\baselineskip}
\pdfbookmark[1]{Ceres Ferdinandea. 1802 Juni}{ceres180206}
\begin{center}
CERES FERDINANDEA.\par
\medskip\hrule\medskip
Monatliche Correspondenz. Band V. S.\ 576—590. 1802 Juni.\par
\medskip\hrule
\end{center}\medskip
—\quad —\quad —

Auch Dr.\ \textsc{Gauß} hat wie bisher alle unsere Beobachtungen dieses Planeten weiter mit den Elementen VII seiner Bahn verglichen. Dieser Vergleich bestätigt, dass sich die Fehler dieser Elemente nur sehr wenig und langsam ändern und die Elemente somit den Ort des Planeten während seiner gesamten diesjährigen Sichtbarkeit mit mehr als ausreichender Genauigkeit angeben. Der Vergleich wird ihm auch künftig sehr zugutekommen, wenn er diese Elemente weiter verfeinert und die Störungsgleichungen einbezieht. Bis zum 19. März hatte Dr.\ \textsc{Gauß} diesen Vergleich im Maiheft auf S.\ 470 [S.\ 211 dieses Bandes] fortgeführt. Hier folgen die Ergebnisse der nächsten zwölf Beobachtungen.
\begin{center}
{\small\setlength{\tabcolsep}{4pt}
\begin{tabular}{r r|r|r|r|r|r}
\multicolumn{2}{c|}{1802}&\shortstack{Mittlere Zeit\\auf dem Seeberg}&\shortstack{Berechnete RA\\der Ceres}&\shortstack{Berechnete\\Deklination\\der Ceres}&\multicolumn{2}{c}{\shortstack{Unterschied\\R—B}}\\
\multicolumn{2}{c|}{}&&&&in RA&in Dekl.\\\hline
März&23&$12^{\mathrm{u}}\quad6^{\mathrm{m}}\ 28^{\mathrm{s}}9$&$182^0\ 20'\ 39''9$&$17^0\ 36'\ 51''4$&$+16''1$&$+23''5$\\
&27&$11\ 47\ 20.3$&$181\ 29\ 21.2$&$17\ 50\ 55.5$&$+20.1$&$+25.9$\\
&28&$11\ 42\ 33.7$&$181\ 16\ 41.4$&$17\ 53\ 52.9$&$+23.7$&$+29.8$\\
&29&$11\ 37\ 48.0$&$181\quad4\quad7.5$&$17\ 56\ 36.3$&$+18.3$&$+32.3$\\
&30&$11\ 33\quad2.5$&$180\ 51\ 40.3$&$17\ 59\quad5.5$&$+16.4$&$+36.5$\\
&31&$11\ 28\ 17.0$&$180\ 39\ 21.1$&$18\quad1\ 20.3$&$+24.1$&$+31.9$\\
April&1&$11\ 23\ 32.4$&$180\ 27\ 10.7$&$18\quad3\ 20.6$&$+25.2$&$+30.8$\\
&2&$11\ 18\ 48.7$&$180\ 15\ 10.1$&$18\quad5\quad6.2$&$+23.6$&$+31.4$\\
&3&$11\ 14\quad6.0$&$180\quad3\ 20.1$&$18\quad6\ 37.4$&$+17.7$&$+27.0$\\
&4&$11\quad9\ 24.1$&$179\ 51\ 41.7$&$18\quad7\ 53.5$&$+10.8$&$+24.1$\\
&5&$11\quad4\ 41.9$&$179\ 40\ 15.6$&$18\quad8\ 54.9$&$+21.5$&$+27.6$\\
&7&$10\ 55\ 21.3$&$179\ 18\quad3.9$&$18\ 10\ 12.5$&$+24.2$&$+25.4$
\end{tabular}}
\end{center}

Die Opposition des Planeten zur Sonne berechnete Dr.\ \textsc{Gauß} aus unseren Beobachtungen vom 15., 16., 17., 18. und 19. März. Er fand, dass sie am 17. März um $4^{\mathrm{u}}\ 21^{\mathrm{m}}\ 10^{\mathrm{s}}$ mittlerer Seeberger Zeit bei einer Länge von $5^{\mathrm{s}}\ 26^0\ 21'\ 27''3$ und einer geozentrischen Breite von $17^0\quad7'\ 59''6$ eintrat.
\par\noindent —\quad —\quad —\hfill [\textsc{von Zach}.]\par

